\section{Model, fidelity and detailed three-dimensional theorem}\label{s:model}
This section fixes the local circuit model, the two noise models and the fidelity and lifetime used throughout. It states Theorem~\ref{thm:detailed}, the detailed form of Theorem~\ref{thm:main}.

Universality supplies a total computable compiler from program codes \(e\) to inputs \(n\) such that \(U(n)\) halts exactly when program \(e\) halts on input \(0\) (Section~\ref{a:certification}). Each construction has nonempty open windows \(I_n\) with rational endpoints. The total computable noiseless encoder depends on \((n,L)\), and the total computable noiseless decoder on \((n,L,t)\). On a halting window the forgetting bound holds for all encoders and decoders on the entire register, for every \(L\), every rate in \(I_n\), and every integer \(t\) at or above the cutoff \(C_{n,\eta}\) in the corresponding theorem. Time counts rounds in Theorems~\ref{thm:exponential} and~\ref{thm:main} and elementary layers in Theorem~\ref{thm:four-dimensional}.

\subsection{Terminology}
We call the registers of the rule before physical compilation \emph{semantic registers}. A site can hold several such registers and the extra wires used to implement them. Table~\ref{tab:terms} gives the recurring terms.

\begin{table}[!tb]
\caption{Terminology used in the construction and the proofs. The source vocabulary (bin, star, seed, witness, certificate, charge, image, catalogue and overlay) is explained in Section~\ref{c:named-sources}.}\label{tab:terms}
\small\centering
\begin{tabular}{@{}>{\raggedright\arraybackslash}p{.23\linewidth}>{\raggedright\arraybackslash}p{.51\linewidth}>{\raggedright\arraybackslash}p{.19\linewidth}@{}}
\toprule Term & Meaning & Defined in\\\midrule
Selector and chart & Rate-selecting classical controller and its locally checked space-time diagram. In six and four dimensions the chart is a static patch. A live letter is a nonbottom symbol. & Sections~\ref{c:chart}, \ref{sec:exponential} and~\ref{sec:four-dimensional}\\\addlinespace[3pt]
Cap, profile and death flag & Finite static computation patch with a diagonal overflow counter. Its per-site profile is \(\mathsf{pr}\). The flag \(\mathsf d\) records whether erasure has been triggered at or has propagated to a site. & Section~\ref{sec:exponential}\\\addlinespace[3pt]
Replica and sampling lane & A replica is one transverse displacement supplying one vote. A sampling lane follows a row by copying a register that samples noise. One such copying interval is called a ride. & Sections~\ref{c:chart} and~\ref{v:parameters}\\\addlinespace[3pt]
Censor & Local statistical test that can replace a computation symbol by \(\bot\). & Section~\ref{sec:exponential}\\\addlinespace[3pt]
Storage-enable bit and cut & In six and four dimensions the storage-enable bit controls conditional reset. A cut removes incoming dependence through that reset. & Sections~\ref{sec:exponential}, \ref{sec:four-dimensional}\\\addlinespace[3pt]
Conditional reset and mixed ancilla & Controlled swap with a fresh maximally mixed ancilla, then reset of that ancilla. In the clean circuit a triggered cut replaces each data qubit by \(I/2\). & Section~\ref{sec:compilation}\\\addlinespace[3pt]
Hidden cover, realization and rate bound & A hidden cover adds letters to projected chart geometry so legality is local. A realization compiles the rule for one noise model. A rate bound is an upper bound on a fault or refresh probability. & Sections~\ref{c:chart} and~\ref{sec:compilation}\\\addlinespace[3pt]
Dirty location & Circuit location touched by a fault or by a wrong controller output. & Section~\ref{d:dispatch}\\\addlinespace[3pt]
Sheet and dispatch & A sheet is one affine branch of a folded code. Dispatch is the selection and execution of a local gate through its ports. & Section~\ref{d:dispatch}, Appendices~\ref{app:f4-dispatch} and~\ref{sec:folded}\\\addlinespace[3pt]
Clock and exclusive dispatcher & The clock is a classical scheduling field. The dispatcher uses exclusive ports to implement the chosen gate. & Section~\ref{d:calendar-section}, Appendices~\ref{app:f4-calendar} and~\ref{app:f4-dispatch}\\
\bottomrule
\end{tabular}
\end{table}
Table~\ref{tab:geometry} lists the coordinate roles. In the six-dimensional construction \(u\) denotes the patch coordinates and \(P\) the four code coordinates. In the four-dimensional construction \(\mathsf C\) is the computation field and \(\mathsf S\) the clock field. The chain coordinate is \(\sigma\), the clock block length is \(P_{\rm cal}\), the simulation period is \(T_{\rm cal}\), and \(k\) is the concatenation depth. The window bound \(w_n\) is on bin height and \(\alpha_4\) is the storage exponent.

\begin{table}[!tb]
\caption{Coordinate roles. The second coordinate \(t-i\) of the three-dimensional computation is a space-time coordinate.}\label{tab:geometry}
\centering\small
\begin{tabular}{@{}l>{\raggedright\arraybackslash}p{.32\linewidth}>{\raggedright\arraybackslash}p{.49\linewidth}@{}}
\toprule Circuit & Computation & Quantum data\\\midrule
6D (A) & Static patch in \((u_x,u_y)\) & Four-dimensional toric code in \((P_1,P_2,P_3,P_4)\)\\\addlinespace[3pt]
4D (B) & Static patch in \((x,y)\) & Chain along \(\sigma\), folded into two internal roles\\\addlinespace[3pt]
3D (C) & Space-time diagram in \((i,t-i)\) & Chain along \(\sigma\)\\\bottomrule
\end{tabular}
\end{table}

The separate four-dimensional circuit of Appendix~\ref{sec:coexistence}
uses the four-sheet geometry of Appendix~\ref{sec:folded} and a separate
closed-code track. The recurring constants are explained in
Appendix~\ref{app:code-estimates}.

\subsection{Local circuits and two noise models}\label{s:local-circuits}
Each construction has a fixed local register and a fixed repeating schedule. We begin with the three-dimensional register convention; the later theorems specify their other lattice dimensions.
Let \(\mathbb T_L^3=(\Z/L\Z)^3\). One physical site carries
\(\mathcal H_x=(\mathbb C^2)^{\otimes N}\), where \(N\) is one fixed
finite integer. It includes every data, controller, sensor, transport
and work qubit. A gate layer is a translated product of gates on
disjoint internal tensor factors. Several such gates can meet the same
physical product site on different factors. Every gate has at most
three input qubits and uses one of finitely many displacement templates, independent of \(n,L,\eps\).

\begin{definition}[Whole-site depolarizing noise, WS]\label{def:noise-w}
Following each gate layer, independently at each physical site, apply
\begin{equation}\label{eq:site-noise}
 \mathcal D_{\eps,x}(\rho)
 =(1-\eps)\rho+\eps\,\frac{I_x}{2^N}\otimes\Tr_x\rho.
\end{equation}
\end{definition}
\begin{definition}[Independent single-qubit depolarizing noise, Q]\label{def:noise-q}
In realization \(\mathrm Q\), independently on every wire after every layer, apply
the one-wire channel \(\rho\mapsto(1-\eps)\rho+\eps I/2\), tensored with the identity on all other wires.
\end{definition}
In both realizations the same
physical rate follows preparation, dephasing, reset, and idle layers.
\(\mathrm{WS}\) and \(\mathrm Q\) are separately compiled fixed circuits; their dimensions,
schedules and physical windows may differ. Gates in a layer are disjoint
on qubits and may share physical sites. All spatial gates have fixed
finite range, which need not be one lattice spacing. The gate list is
\[
 \{X,H,\mathrm{CNOT},\mathrm{Toffoli},R_0,D_Z,I\}.
\]
The unitary subset contains the computationally universal pair Toffoli and Hadamard \cite{Aharonov2003,Shi2003} (universality here is the real-amplitude simulation of complex computation with one extra qubit, in the sense of \cite[Definition~3 and Theorem~2]{Aharonov2003}). Here \(R_0(\rho)=|0\rangle\langle0|\Tr\rho\) and \(D_Z\) is complete
computational-basis dephasing; a measurement is a \(D_Z\) whose outcome
stays on a basis-valued wire.
\Needspace{7\baselineskip}
\begin{definition}[Round channel]\label{def:round}
Repeating the fixed circuit defines a round channel \(\Lambda_{L,\eps}\).
One application is one round. If the period is \(J\) elementary layers,
then at aligned times \mbox{\(t_{\rm layer}=Jt_{\rm round}\)}.
For Theorem~\ref{thm:four-dimensional}, including Theorem~\ref{f4:detailed},
bound~\eqref{f4:lifetime-upper}, Table~\ref{tab:results} and
Corollary~\ref{cor:window-undecidable}, fidelity and lifetime are measured
in elementary layers. In this convention use \(\Lambda^{(t)}_{L,\eps}\),
the channel of the first \(t\) noisy layers starting at a round boundary,
in Definitions~\ref{def:fidelity} and~\ref{def:lifetime}. Convert round
quantities in the proof at aligned times using the period.
Unless elementary layers are specified, \(t\) counts rounds.
\end{definition}

\subsection{Storage fidelity}
CPTP means completely positive and trace preserving.
\begin{definition}[Entanglement fidelity]\label{def:fidelity}
Let \(\beta=|\beta\rangle\langle\beta|\), with
\(|\beta\rangle=(|00\rangle+|11\rangle)/\sqrt2\). The reference qubit
is never acted upon. For CPTP encoding and decoding maps \(\cE,\cD\), put
\begin{equation}\label{eq:fidelity}
 F(\cE,\cD;L,t,\eps)=
 \langle\beta|[\mathrm{id}\otimes
 (\cD\Lambda_{L,\eps}^t\cE)](\beta)|\beta\rangle,
 \qquad F^*_{L,\eps}(t)=\sup_{\cE,\cD}F(\cE,\cD;L,t,\eps).
\end{equation}
This is entanglement fidelity for the maximally mixed logical qubit.
\end{definition}
The supremum allows arbitrary encodings into the entire physical
algebra, including work and controller registers, with no persistent
outside quantum memory. \begin{definition}[Memory lifetime, or storage time]\label{def:lifetime}
Let
\[
 \tau_{L,\eps}=\sup\{t\in\N:F^*_{L,\eps}(t)\ge7/8\}.
\]
\end{definition}
Optimal fidelity is nonincreasing in \(t\). A later decoder composed with the intervening channel is a decoder at the earlier time. The passing layer times
therefore form an initial segment, and the passing round times are
precisely the aligned times in that segment. For a period of \(J\)
layers, the optimized integer lifetimes satisfy
\(\tau_{\rm round}=\lfloor\tau_{\rm layer}/J\rfloor\) when finite,
and are both infinite otherwise. Our explicit decoder can depend
on \(t\), as well as \(n,L\), and remains independent of the rate.


\begin{definition}[Storing and forgetting windows]\label{def:windows}
A \emph{rate window} is a nonempty open interval with rational endpoints. For a specified encoder family, decoder family, horizon and size threshold, a \emph{storing window} is one on which the stated fidelity guarantee holds at every rate and every indicated size and time. A \emph{forgetting window} is one on which, for every \(\eta>0\), the optimal fidelity is at most \(1/2+\eta\) after a finite time uniform in size and rate. A \emph{scalable-memory rate} is one with \(\sup_L\tau_{L,\eps}=\infty\). A \emph{scalable-memory window} consists entirely of such rates. Nonhalting and halting windows are indexed by \(n\) with \(U(n)\uparrow\) and \(U(n)\downarrow\), respectively. The theorems give storing windows with fidelity \(15/16\); certification tests the weaker bound \(7/8\). The limiting benchmark \(1/2\) is the largest entanglement fidelity attainable by a measure-and-prepare qubit channel.
\end{definition}

\begin{theorem}[Detailed form of Theorem~\ref{thm:main}]\label{thm:detailed}
For each of \(\mathrm{WS}\) and \(\mathrm Q\) there are a computably specified,
translation-invariant, time-periodic circuit as above, a computable sequence
of disjoint nonempty rational open intervals
\(I_n\subset(0,1)\), a rational \(c>0\), and computable thresholds
\(L_0(n)\), with \(c\) independent of \(n\), and the following properties.
\begin{enumerate}
\item There are total computable maps \(\cE_{n,L}\) and
\(\cD_{n,L,t}\), independent of \(\eps\) and of whether \(U(n)\)
halts. If \(U(n)\uparrow\), then for every \(\eps\in I_n\),
\(L\ge L_0(n)\), and every integer \(0\le t\le\lfloor cL\rfloor\),
\(F(\cE_{n,L},\cD_{n,L,t};L,t,\eps)\ge15/16\).
\item If \(U(n)\downarrow\), then for every \(\eta>0\) there is
\(C_{n,\eta}<\infty\) such that
\(F^*_{L,\eps}(t)\le1/2+\eta\) for every \(L\),
\(\eps\in I_n\), and \(t\ge C_{n,\eta}\).
\item There are computable constants \(A,a>0\), \(b\in\N\), \(C_0\ge2\)
and \(L_1\), all independent of \(n\), such that for every \(n\), every
\(\eps\in I_n\) and every \(L\ge L_1\), with \(s_L=201L+C_0\),
\[
 \Lambda_{L,\eps}^{s_L}=(1-r_L)\Phi^{\rm EB}_{L,\eps}+r_L\mathcal B_{L,\eps},
 \qquad r_L\le q_L\le AL^be^{-aL},
\]
where \(\Phi^{\rm EB}_{L,\eps}\) is entanglement breaking (EB) and
\(\mathcal B_{L,\eps}\) is a channel. In particular
\mbox{\(\|\Lambda_{L,\eps}^{s_L}-\Phi^{\rm EB}_{L,\eps}\|_\diamond\le2q_L\)} and
\(F^*_{L,\eps}(t)\le1/2+q_L/2\) for every \(t\ge s_L\).
\item The predicate requiring the explicit fidelity to be at least
\(7/8\) for every integer \(L\ge L_0(n)\), integer
\(0\le t\le\lfloor cL\rfloor\), and rational \(\eps\in I_n\)
holds exactly when \(U(n)\uparrow\) and is \(\Pi^0_1\)-complete.
Every finite test is an exact rational comparison.
\end{enumerate}
\end{theorem}

If \(U(n)\) never halts, these bounds give
\(\lfloor cL\rfloor\le\tau_{L,\eps}<201L+C_0\) for every \(\eps\in I_n\)
and every \(L\ge L_C(n):=\max\{L_0(n),L_1,L_2\}\). Here \(L_2\) is a
computable size beyond which \(AL^be^{-aL}<3/4\). One can search beyond
\(b/a\), where this bound decreases, until a certified upper estimate
is strictly below \(3/4\). Thus
\(\tau_{L,\eps}=\Theta(L)\), uniformly across the window. On halting windows the lifetime is
uniformly bounded in \(L\). These halting-side times are finite but not
computable from \(n\): for each fixed rational \(0<\eta<3/8\), no total
computable function of \(n\) can serve as \(C_{n,\eta}\) in part (ii)
(Corollary~\ref{a:time}). Computable windows and their rational midpoints
already suffice for the undecidability assertion. The scope is as in Section~\ref{intro:scope}.

For fixed \(L\) and \(\eps>0\), complete replacement of every physical
factor at the last layer of a round has probability
\(p_{L,\eps}=\eps^{L^3}\) in \(\mathrm{WS}\) and
\(p_{L,\eps}=\eps^{NL^3}\) in \(\mathrm Q\). If it has occurred, the output is
independent of the input. Thus
\begin{equation}\label{eq:finite-reset}
 F^*_{L,\eps}(t)\le\frac14+\frac34(1-p_{L,\eps})^t.
\end{equation}
This also controls the finitely many small volumes omitted from
asymptotic estimates.

Sections~\ref{sec:controller}--\ref{a:certification} prove Theorem~\ref{thm:detailed}. Section~\ref{sec:controller} constructs the classical selector and Section~\ref{sec:compilation} compiles it and fixes the constants. Section~\ref{sec:quantum} proves (i), Section~\ref{sec:forgetting} proves (ii), Section~\ref{sec:ceiling} proves (iii), and Section~\ref{a:certification} proves (iv).
